\section{Bounded Depth Implies Bounded Length}
\label{sec:search}

The last bound turns a bound on the depth into a bound on the length. Like the depth bound,
it holds for every reader: it uses neither the admission cone nor any property of
non-structural subtype entailment.

\begin{theorem}[Bounded-depth search]\label{thm:search}
Let $(\mathcal{M},\Adm,\Rej)$ be any reader with $N$ states over a finite alphabet. If it has a
nonempty hole of depth at most $d$, then it has a nonempty hole of length at most
$\holeBound(N,d)$, where $\holeBound$ is the elementary function given in
\eqref{eq:holebound} below. The replacement can be chosen of depth at most $d$.
The same statement holds for any lettered monitor with $N$ states
(defined in \cref{sec:lettered}), even over an infinite alphabet.
\end{theorem}

The proof first bounds the number of distinct letters at every positive
level, then shortens a hole level by level along its hierarchy. Shortening each block of a
level independently fails, because comparisons relate \emph{different} blocks through equal
suffixes. We therefore shorten all blocks of a level \emph{simultaneously}, preserving a
finite summary of the whole tuple. As \cref{ex:loop-deletion} shows,
preserving the final reader state alone does not preserve holes.

\subsection{Lettered Monitors and Their Derivatives}\label{sec:lettered}

To describe each level of the hierarchy as a hole problem in its own right, we generalize
readers to \emph{lettered monitors}. A lettered monitor $\mathcal{M}$ over an alphabet
$\Gamma$, which may be infinite, consists of a finite set $\mathcal Q$ of states with a start state $\zeta_0$, a transition
function $\delta : \mathcal Q \times \Gamma \to \mathcal Q$, a \emph{final target} predicate
$\Tfin \subseteq \mathcal Q \times \Gamma$, and two admission relations $\Lambda$ and $\Lfin$ on
\emph{labels}, the second for comparisons that end at the final cut. Reading a word $v$
assigns the state $\zeta_{i+1} = \delta(\zeta_i, v[i])$ to cut $i+1$. The label of cut $0$ is a
symbol $\triangleright$, and the label of cut $i > 0$ is the pair $(\zeta_{i-1}, v[i-1])$ of the
state at cut $i-1$ and the letter read from it. A nonempty word $v$ is a \emph{hole} of
$\mathcal{M}$ if its last label $(\zeta,a)$ satisfies $\Tfin(\zeta,a)$, no terminal comparison
$(s,\len{v})$ satisfies $\Lfin$, and no other comparison $(s,e)$ satisfies $\Lambda$, each
relation being applied to the labels of the two cuts. A reader $(\mathcal{M},\Adm,\Rej)$ determines the
lettered monitor with states $Q$, the same transition $\delta$, $\Tfin(\zeta,a)$ iff $\delta(\zeta,a) \in \Rej$,
and $\Lambda = \Lfin$ given by $\Adm$ on the states that the labels denote: $\triangleright$
denotes $q_0$ and $(\zeta,a)$ denotes $\delta(\zeta,a)$. The reader and this lettered monitor have the same nonempty holes.
As with rejection, the final target only marks a candidate: a word whose last label
satisfies $\Tfin$ is a hole only if no comparison is admitted.
The predicate $\Tfin$ and the relations take labels, and not only the states
they denote, because the monitors of higher levels must evaluate them at positions
\emph{inside} blocks, which requires the letters. Block alphabets $(\Gamma\setminus\{\ell\})^*$
are generally infinite, even for finite $\Gamma$, so lettered monitors allow infinite
alphabets. The compression argument never uses finiteness of an alphabet: it works with
the finitely many letters that one hole uses and with finitely many summary values.

Fix a marker $\ell \in \Gamma$, a block $X \in (\Gamma \setminus \{\ell\})^*$, and a state
$\zeta \in \mathcal Q \sqcup \{\star\}$. The \emph{shifted block} of $X$ read from $\zeta \in \mathcal Q$ is the word
$\ell X$ read from $\zeta$; from the new symbol $\star$, it is $X$ read from $\zeta_0$. It has a label
$\mathit{lab}_p(\zeta,X)$ at each position $p = 0, \dots, \len{X}$: for $\zeta \in \mathcal Q$, the label of the
cut after $\ell\,\wpre{X}{p}$, so that position $0$ carries the \emph{marker label}
$(\zeta,\ell)$; for $\zeta = \star$, the label $\triangleright$ at position $0$ and the label of the cut
after $\wpre{X}{p}$ at positions $p \ge 1$. The \emph{exit state} $\delta_\ell(\zeta,X)$ is the
state reached at the end of the shifted block.

The derivative $\mathcal M_\ell$ defined next reads the block word. Each of its
letters represents an entire block, so the label of a derivative cut keeps both
the block and its entry state, the state from which the shifted block is read;
this lets the derivative recover every original label inside the block. With
these labels, $\mathcal M_\ell$ checks three kinds of obligations of the expanded
word. Its admission $\Lambda_\ell$ compares two positions in different blocks
that lie at the same distance from the ends of their blocks, with equal letters
after them. The second clause of its terminal admission $\Lfin_\ell$ compares a
position anywhere in an earlier block with the final cut, which carries the label
of the closing marker. Its target $\Tfin_\ell$ checks the original target together
with all terminal comparisons that start in the last block.

\begin{definition}[Derivative]\label{def:derivative}
Let $\mathcal{M}$ be a lettered monitor over $\Gamma$ with states $\mathcal Q$, and $\ell \in \Gamma$.
The \emph{derivative} $\mathcal{M}_\ell$ is the lettered monitor over the block alphabet
$(\Gamma\setminus\{\ell\})^*$ with states $\mathcal Q \sqcup \{\star\}$, start state $\star$, transition
function $\delta_\ell$, and the following relations on labels $(\zeta,X)$ and $(\zeta',X')$:
\begin{itemize}[leftmargin=*]
\item $\Lambda_\ell((\zeta,X),(\zeta',X'))$ holds if
  $\Lambda\bigl(\mathit{lab}_{p}(\zeta,X),\, \mathit{lab}_{p'}(\zeta',X')\bigr)$ for some
  positions $p\le\len{X}$ and $p'\le\len{X'}$ from which the two blocks have a common
  suffix, $\wsuf{X}{p}=\wsuf{X'}{p'}$;
\item $\Lfin_\ell((\zeta,X),(\zeta',X'))$ holds if $\Lambda_\ell((\zeta,X),(\zeta',X'))$ holds, or if
  $\Lfin\bigl(\mathit{lab}_p(\zeta,X), (\delta_\ell(\zeta',X'),\ell)\bigr)$ for some
  $p \le \len{X}$;
\item $\Tfin_\ell(\zeta,X)$ holds if $\Tfin(\delta_\ell(\zeta,X),\ell)$ holds and
  $\Lfin\bigl(\mathit{lab}_p(\zeta,X), (\delta_\ell(\zeta,X),\ell)\bigr)$ fails for every
  $p \le \len{X}$.
\end{itemize}
Every relation involving the label $\triangleright$ of $\mathcal{M}_\ell$ is false.
\end{definition}

The derivative's cut zero is an artificial boundary. Original comparisons
beginning at cut zero are represented inside the first shifted block,
whose entry state is $\star$; they are still tested through the label
$(\star,X_0)$. Setting relations at the derivative's new start label to
false therefore discards no original comparison.

For a concrete instance, use the reader of \cref{ex:loop-deletion}
on $0101$ and marker $1$. Its block word is $[0][0]$. The derivative
starts at $\star$ and reads $[0]$ to state $a$; from $a$, the shifted
block $10$ ends at $c$. The table shows which derivative label represents
each original cut.
\begin{center}
\small
\begin{tabular}{@{}lccccc@{}}
\toprule
original cut & $0$ & $1$ & $2$ & $3$ & $4$\\
state & $q$ & $a$ & $b$ & $c$ & $r$\\
original label & $\triangleright$ & $(q,0)$ & $(a,1)$ & $(b,0)$ & $(c,1)$\\
\midrule
derivative label & \multicolumn{2}{c}{$(\star,[0])$} & \multicolumn{2}{c}{$(a,[0])$} & final\\
\cmidrule(lr){2-3}\cmidrule(lr){4-5}
position $p$ in its block & $0$ & $1$ & $0$ & $1$ & \\
\bottomrule
\end{tabular}
\end{center}
The two derivative labels sit at the derivative cuts $1$ and $2$, and each
represents the original cuts of one shifted block: the first block is read from
the start, and the second begins with the marker label $(a,1)$. The final original
cut carries the label $(c,1)$ of the closing marker, read from the final derivative
state $c$. For the two derivative labels, the common suffix $0$ selects position
zero in both blocks. Their original labels denote $q$ and $\delta(a,1)=b$, so
$\Lambda_1$ holds. The derivative comparison $(1,2)$ is terminal and its $\Lfin_1$
relation therefore holds as well. It detects the original nonterminal comparison
$(0,2)$, including its start at the original cut zero.

\begin{lemma}[Holes under desubstitution]\label{lem:der-holes}
Let $v$ be a nonempty word of blocks that avoid $\ell$, and let $u$ be its expansion. Then
$u$ is a hole of $\mathcal{M}$ iff $v$ is a hole of $\mathcal{M}_\ell$. Consequently, if
$\mathcal{M}_0$ is a reader with $N$ states and $\mathcal{M}_{j+1} = (\mathcal{M}_j)_{\ell_j}$
along the canonical hierarchy of $w$, then $w$ is a hole of $\mathcal{M}_0$ iff $u_j$ is a
hole of $\mathcal{M}_j$, and $\mathcal{M}_j$ has $N + j$ states.
\end{lemma}

\begin{proof}
Write $v=[X_0]\cdots[X_{h-1}]$, so
$u=X_0\ell X_1\ell\cdots X_{h-1}\ell$, and let
$\zeta'_0=\star,\zeta'_1,\ldots,\zeta'_h$ be the derivative run.
Let $a_i$ be the cut where $X_i$ begins and $b_i=a_i+|X_i|$ the
position of its closing marker. Partition all nonfinal cuts of $u$ into
the intervals $[a_i,b_i]$. Induction on $i$ gives the label
$\mathit{lab}_p(\zeta'_i,X_i)$ at cut $a_i+p$, for $0\le p\le|X_i|$;
the final cut has label $(\zeta'_h,\ell)$. The label of cut $i+1$ of $v$
is $(\zeta'_i,X_i)$, so it determines every label in this block.

Suppose first that $u$ is a hole. Its final target holds, and none of
its terminal comparisons starting in the last block is admitted. These
are exactly the two requirements of $\Tfin_\ell(\zeta'_{h-1},X_{h-1})$.
Consider a comparison $(s',e')$ of $v$. If $s'=0$, its relations are false
by definition. Otherwise, $1\le s'<e'\le h$. If $\Lambda_\ell$ held at its
labels, it would supply positions $p$ in $X_{s'-1}$ and $p'$ in $X_{e'-1}$
with $\wsuf{X_{s'-1}}{p}=\wsuf{X_{e'-1}}{p'}$.
The cuts $s=a_{s'-1}+p$ and $e=a_{e'-1}+p'$ have this common
suffix before their markers. By \cref{lem:alignment}, $(s,e)$ is
a comparison of $u$; it is nonterminal even when $e'=h$, since
$e\le b_{h-1}<|u|$. Its labels would satisfy $\Lambda$, a contradiction.
This excludes $\Lambda_\ell$ both at nonterminal comparisons of $v$
and as the first disjunct of $\Lfin_\ell$ at terminal ones.
The second disjunct, possible only for $e'=h$, gives an admitted terminal
comparison $(a_{s'-1}+p,|u|)$ of $u$ for some position $p$, also a
contradiction. Thus $v$ is a hole.

Conversely, suppose $v$ is a hole. Its target implies the final target
of $u$ and excludes every admitted terminal comparison starting in the
last block. Let $(s,e)$ be any other comparison of $u$. If $e<|u|$,
\cref{lem:alignment} locates $s,e$ in blocks $s'-1,e'-1$ with $s'<e'$,
a common suffix after them, and a comparison $(s',e')$ of $v$.
Admission by $\Lambda$ would therefore give $\Lambda_\ell$ at its labels.
For $e'<h$ this contradicts the nonterminal hole condition of $v$;
for $e'=h$ it contradicts the first disjunct of its terminal relation
$\Lfin_\ell$. If $e=|u|$, its start is in an earlier block $s'-1$
with $s'<h$. Then $(s',h)$ is a terminal comparison of $v$, and admission
by $\Lfin$ would give the second disjunct of $\Lfin_\ell$.
Again this is impossible. All comparisons of $u$ are therefore unadmitted,
so $u$ is a hole. Iterating this equivalence proves the hierarchy assertion;
each derivative adds exactly one fresh state.
\end{proof}



\subsection{Isolated Exchange and Bounded Dictionaries}
\label{sec:dictionary}

Shortening (\cref{sec:shortening}) compresses all distinct blocks of a level at
once, and its length bound grows quickly with the number of these blocks. We
therefore need a bound on the number of different blocks at each positive level,
not just on their lengths. We obtain it by merging letters that have the same
effect on the derived monitor. Such a merge must preserve the comparisons of the
level word as well as its transitions, targets, and admissions. We therefore
protect the marker and the letters at symbolic start positions, and merge only
letters whose unequal comparison partners are all protected. The behavior of such
a letter can then be recorded in finite tables, its exchange type, whose size does
not depend on the dictionary. The other letters are counted through the boundary
pairs defined below, and these counts propagate from one level to the next.

Fix a nonempty hole $w$ of depth $\Delta\le d$, its hierarchy $u_i$, and
its derived monitors $\mathcal M_i$. For a word $v$, let
$\operatorname{alph}(v)$ be its set of letters. Write
$\mathcal A_i=\operatorname{alph}(u_i)$ for the level-$i$
\emph{dictionary}, and
$g_i=|\mathcal A_i|$. These are letter dictionaries, distinct from the
state supports $\Supp_i$ of \cref{sec:depth}. A comparison at level $i$
is a comparison of $u_i$ in the sense of \cref{def:comparison}.
It is \emph{left maximal} if $1\le s<e\le |u_i|$ and its predecessor
letters $u_i[s-1]$ and $u_i[e-1]$ differ, so that it cannot be extended one
letter to the left. Comparisons that begin at cut zero are excluded: they
have no predecessor letter, and their admissions involve the start label,
which the exchange types below record separately. The \emph{boundary pairs}
of level $i$ are the ordered letter pairs
\[
 \operatorname{BP}_i=\{(u_i[s-1],u_i[e-1]):
       (s,e)\text{ is left maximal},\ e<|u_i|\},
\]
counted once regardless of how many comparisons produce them; put
$b_i=|\operatorname{BP}_i|$. A left-maximal comparison that ends at the
final cut has the marker $\ell_i$ as its second predecessor letter, since
$u_i$ ends with $\ell_i$. Its pair, an endpoint pair, is needed only in the
proofs (\cref{app:search}).

The positions of $u_i$ with start cores form a prefix, and their number
$s_i\le i$ is the \emph{start count} of level $i$ (\cref{lem:cores}). In the
derived monitor, ordinary positions are read from the embedded copy of the
original state set $Q$, and start positions from symbolic states. Let $Z_i$
be the union of that copy of $Q$ and the symbolic states used at start
positions; then $|Z_i|\le N+s_i$. The \emph{protected letters} of level $i$ are
\[
 \Pi_i=\{\ell_i\}\cup\operatorname{alph}(u_i[:s_i]),
 \qquad |\Pi_i|\le s_i+1.
\]
A letter $a\in\mathcal A_i\setminus\Pi_i$ is \emph{isolated} if no
boundary pair has $a$ as one component and a letter outside $\Pi_i$ as the
other. All unequal comparison partners of an isolated letter are therefore
protected, including its partner in terminal comparisons, the marker.

\begin{example}[Protected and isolated letters]\label{ex:isolated}
At level $1$ of \cref{ex:hierarchy}, the word $u_1=\alpha_1\beta_1\gamma_1\alpha_1$
has cores $\star_1,1,4,5$ and marker $\ell_1=\alpha_1$. Its start count is
$s_1=1$ and its start prefix is $\alpha_1$, so $\Pi_1=\{\alpha_1\}$. The only
comparison of $u_1$ that ends before the final cut is $(0,3)$, which begins at
cut zero. Hence $\operatorname{BP}_1=\varnothing$, and $\beta_1$ and $\gamma_1$
are isolated. Their only unequal comparison partner is the marker, through the
terminal comparisons $(2,4)$ and $(3,4)$. Merging $\beta_1$ and $\gamma_1$ would
therefore preserve the comparisons of $u_1$; it would also require equal exchange
types, defined next, which depend on the monitor.
\end{example}

For an isolated letter $a$, its \emph{exchange type relative to $\Pi_i$}
records the following finite tables in $\mathcal M_i$:
\begin{itemize}[leftmargin=*]
\item the transition $q\mapsto\delta_i(q,a)$ and target
  $\Tfin_i(q,a)$ for every ordinary state $q\in Q$;
\item for each $R\in\{\Lambda_i,\Lfin_i\}$, the self table
  $R((q,a),(q',a))$ for $q,q'\in Q$;
\item for each such $R$, both $R((q,a),(\zeta,\pi))$ and
  $R((\zeta,\pi),(q,a))$ for $q\in Q$, $\zeta\in Z_i$, $\pi\in\Pi_i$;
\item the two initial-label rows $R(\triangleright,(q,a))$ for $q\in Q$.
\end{itemize}
At positive levels the last rows are false. Original cut-zero obligations
are still recorded in the protected rows at symbolic entry states
(rows $R((q,a),(\zeta,\pi))$ and their reverse orientations, with $\zeta$ symbolic
and $\pi$ in the start prefix).
Ordinary states transition to ordinary states, so the transition table has
$N^N$ possibilities, rather than $(N+i)^N$.

Put
\begin{equation}\label{eq:dictionary-constants}
\begin{aligned}
 \TyBound(N,d)&=N^N\,2^{2N^2+4N(N+d)(d+1)+3N},\\
 \AlBound(N,d)&=d+1+\TyBound(N,d),&
 \DictBound(N,d)&=2+\AlBound(N,d)3^d.
\end{aligned}
\end{equation}
Here $\TyBound$ bounds the number of exchange types, $\AlBound$ adds the
protected letters, and $\DictBound$ will bound every positive dictionary.
Shortening a tuple of $h$ blocks produces blocks shorter than
$\mathit{col}(M,h)$, the column count of \eqref{eq:col} for a monitor with $M$
states, which for fixed $M$ is exponential in $h^2$. Bounding every positive
dictionary by $\DictBound(N,d)$ in advance makes each of at most $d$ expansion
steps multiply lengths by at most the fixed factor
$K=\mathit{col}(N+d+1,\DictBound(N,d))$. The exchange argument achieves such a bound
by merging only isolated letters: their types need rows only against the
at most $d+1$ protected letters, so $\TyBound(N,d)$ is independent of dictionary
sizes. Extending this rule to non-isolated letters would require protecting
their unequal comparison partners as well; their number would enter the
type-count exponent and could compound from level to level.

\begin{lemma}[Exchange types]\label{lem:exchange-types}
At every positive level $i<\Delta\le d$, there are at most $\TyBound(N,d)$
exchange types relative to $\Pi_i$.
\end{lemma}

\begin{lemma}[Isolated exchange]\label{lem:isolated-exchange}
Let $1\le i<\Delta$, and let distinct isolated letters $a,a'$ of $u_i$
have the same exchange type. Replace every $a$ by $a'$ at that level
and expand using the markers of the original hierarchy below level $i$. The
result is a hole of the same canonical depth, with every start count unchanged
and no dictionary larger. Its level-$i$ dictionary is strictly smaller. Choosing
the direction toward the letter with shorter level-zero expansion also makes the
word no longer.
\end{lemma}

Call a hole \emph{dictionary-normal} if, at every $1\le i<\Delta$, distinct isolated
letters have distinct exchange types, computed in that hole's own hierarchy.
\begin{lemma}[Dictionary normal form]\label{lem:dictionary-normal-form}
Every nonempty hole has a dictionary-normal hole of the same depth, no greater
length, the same start counts, and no larger dictionary at any level.
\end{lemma}
\begin{lemma}[Dictionary counting and propagation]\label{lem:dictionary-count}
For a dictionary-normal hole of depth $\Delta\le d$, and $1\le i<\Delta$,
$g_i\le \AlBound(N,d)+2b_i$.
For every nonempty word and every level $i$,
\begin{equation}\label{eq:boundary-propagation}
 b_i\le b_{i+1}+g_{i+1}-1.
\end{equation}
\end{lemma}
\begin{lemma}[Bounded dictionaries]\label{lem:dictionary-bound}
For a dictionary-normal hole of depth $\Delta\le d$, at $1\le i<\Delta$,
\begin{equation}\label{eq:dictionary-recurrence}
 2b_i\le(\AlBound(N,d)-1)(3^{\Delta-1-i}-1),\qquad
 g_i\le1+(\AlBound(N,d)-1)3^{\Delta-1-i}.
\end{equation}
Consequently every nonempty hole of depth at most $d$ has a hole of the
same depth and no greater length with $g_i\le \DictBound(N,d)$ for every
$1\le i\le \Delta$. No bound on $g_0$ is required.
\end{lemma}

Full proofs, including preservation of start counts under the induced
level maps, appear in \cref{app:search}. The counting idea is simple:
dictionary normality bounds isolated letters by their types, while each remaining
unprotected letter is a component of a boundary pair. This adds at most
$2b_i$ letters. Those letters need not enter the protected set
inside the type-count exponent: an isolated letter never compares with
an unprotected unequal partner. Propagation therefore gives a linear
recurrence with factor three, while the type bound remains fixed.

\subsection{Shortening}\label{sec:shortening}

Fix the dictionary-normal hole and its dictionaries $\mathcal A_j$. Throughout
shortening, the original hierarchy means the hierarchy of this hole, which stays
unchanged. Put $g=\DictBound(N,d)$ and
$K=\mathit{col}(N+d+1,g)$, using the column count of \eqref{eq:col}.
We shorten top-down.
First the top level is shortened; then the blocks that form its
letters; then the blocks that form the letters of those blocks; and so on down to level~$0$.

\paragraph{The Top Level}
At the first unary level $\Delta\le d$, write $u_\Delta = z^L$. Every pair of its cuts is a
comparison, and its states follow the orbit of $\delta(\cdot,z)$. Deleting the letters between
two nonfinal cuts with equal states creates no new kind of label pair, so a unary hole of
$\mathcal{M}_\Delta$ can be shortened to at most $N + \Delta$ letters, the number of states
(details in \cref{app:search}).

\begin{figure}[tbp]
\centering
\begin{tikzpicture}[>=Stealth,font=\small,
  stage/.style={draw,rounded corners=2pt,align=center,text width=32mm,
    minimum height=14mm,inner sep=3pt}]
  \node[stage] (unary) at (0,0)
    {unary hole at level $\Delta$\\$z^L$, $1\le L\le N+\Delta$};
  \node[stage] (compressed) at (4.45,0)
    {renamed letter $\rho_\Delta(z)$,\\same monitor values\\$\rho_\Delta(z)^L$};
  \node[stage] (expanded) at (8.9,0)
    {level-$0$ hole $v$\\$|v|\le L\,K^\Delta$\\$\le\holeBound(N,d)$};
  \draw[->] (unary) -- (compressed)
    node[midway,above=8mm,font=\footnotesize] {compress};
  \draw[->] (compressed) -- (expanded)
    node[midway,above=8mm,font=\footnotesize] {expand};
  \node[below=4mm of compressed,align=center,font=\footnotesize]
    {one distinct letter to preserve at level $\Delta$;\\
     shorten its blocks and their letters recursively down to level $0$};
\end{tikzpicture}
\caption{Unary top level and compression. First shorten the unary power by
deleting repeated-state loops, giving $L\le N+\Delta$. Compression then renames
its single letter by the injective renaming $\rho_\Delta$ of
\cref{lem:simultaneous-compression}, under which the new level-$\Delta$ monitor
agrees with the original one: it has the same transition, target, and admission
values on the renamed letter. This preserves the run, target, and admissions of the power.
With the dictionary-normal hierarchy unchanged and
$K=\mathit{col}(N+d+1,\DictBound(N,d))$, supported block compression bounds
that letter's expansion by $K^\Delta$, giving
$|v|\le LK^\Delta\le\holeBound(N,d)$. At $\Delta=0$ the expansion is the identity.}
\Description{Three boxes show a unary hole with at most N plus Delta copies of one
letter, an injective renaming preserving monitor agreement, and expansion to a
level-zero hole of length at most L times the Delta-th power of the fixed column bound.}
\label{fig:unary-compression}
\end{figure}


\paragraph{The Block Levels}
At level $j+1$, select distinct original blocks $X_1,\dots,X_h$ from
$\mathcal A_{j+1}$. They are words over $\mathcal A_j$ avoiding the
original marker $\ell_j$. The \emph{summary} of the tuple
$(X_1,\dots,X_h)$ with respect to $\mathcal{M}_j$ and $\ell_j$ records:
\begin{itemize}[leftmargin=*]
\item for each block and each state $\zeta \in \mathcal Q \sqcup \{\star\}$, the exit state of its shifted
  block read from~$\zeta$;
\item for each block, each such $\zeta$, and each state $\zeta''\in\mathcal Q$, whether some label of the shifted
  block read from $\zeta$ is related by $\Lfin$ to the final label $(\zeta'',\ell_j)$;
\item which blocks are equal;
\item for each ordered pair of blocks and each pair of states $\zeta, \zeta'$, whether
  $\Lambda_{\ell_j}$ relates the two labels.
\end{itemize}
The summary determines every value of $\mathcal{M}_{j+1}$ on the letters $X_1,\dots,X_h$,
together with their equalities. Storing all equality bits preserves
disequality as well: distinct old letters get distinct replacements, so
the replacement is injective. It can be computed by a deterministic finite automaton
reading the tuple as a single word of \emph{columns}: align the blocks at their right ends,
pad them on the left, and read the $h$-tuples of optional letters from left to right
(\cref{fig:columns}). Right alignment is what makes this possible: $\Lambda_{\ell_j}$
compares positions at equal distance from the ends of two blocks with equal letters after
them, so appending a column either relates the two new last positions, or extends the
common suffix of a pair related earlier by an equal letter. A column word is
\emph{valid} if each of its tracks consists of absent entries followed by letters, as
padding on the left requires. If the column word has at least as many columns as there
are summary values, two cuts have the same summary. Deleting the columns between them
preserves the summary. It also preserves validity without any extra state in the
automaton: deleting columns keeps every track of the form ``absent entries, then
letters'', and each column itself records which tracks are present. Repeating this, we
obtain blocks $\widetilde X_1,\dots,\widetilde X_h$ of length less than the number
$\mathit{col}(N+j,h)$ of summary values, with the same summary. Substituting
$\widetilde X_i$ for $X_i$ yields a word that is again a hole of $\mathcal{M}_{j+1}$.

\begin{figure}[tbp]
\centering
\begin{tikzpicture}[x=9mm,y=7mm,>=Stealth,
    cell/.style={draw, minimum width=8mm, minimum height=6mm, inner sep=1pt, font=\small}]
  \node[anchor=east,font=\small] at (-0.7,0) {$X_1$};
  \node[anchor=east,font=\small] at (-0.7,-1) {$X_2$};
  \node[cell] at (0,0) {$0$};
  \node[cell] at (1,0) {$1$};
  \node[cell] at (2,0) {$0$};
  \node[cell] at (3,0) {$2$};
  \node[cell,fill=black!10] at (0,-1) {};
  \node[cell,fill=black!10] at (1,-1) {};
  \node[cell] at (2,-1) {$1$};
  \node[cell] at (3,-1) {$2$};
  \draw[thick,rounded corners=2pt] (1.52,0.5) rectangle (2.48,-1.5);
  \draw[->] (-0.45,-2.0) -- (3.45,-2.0) node[right,font=\footnotesize] {columns read left to right};
  \node[font=\footnotesize,align=left,anchor=west] at (4.2,-0.5)
    {the boxed positions have the same\\distance from their block ends\\and the same remaining suffix;\\$\Lambda_\ell$ tests their labels};
\end{tikzpicture}
\caption{The column encoding of a tuple of blocks, here $X_1 = 0102$ and $X_2 = 12$. The
blocks are aligned at their right ends and padded on the left (shaded cells). A deterministic
automaton reads the columns from left to right and maintains the summary of the tuple: the
exit states of the shifted blocks, their tests against the final label, the equalities between blocks, and the
bits of $\Lambda_\ell$. Deleting the columns between two equal summaries shortens all blocks
at once and preserves the summary (\cref{app:search}).}
\Description{Two blocks aligned at their right ends, the shorter one padded on the left, with one column highlighted.}
\label{fig:columns}
\end{figure}

The recursion keeps a set $\mathcal E_{j+1}\subseteq\mathcal A_{j+1}$ of distinct
original letters, its \emph{letter support}. After shortening its blocks to
subsequences $\tau(X)$, form the lower letter support
\begin{equation}\label{eq:descent-support}
 \mathcal E_j=\{\ell_j\}\cup\bigcup_{X\in\mathcal E_{j+1}}\operatorname{alph}(\tau(X))
 \ \subseteq\ \mathcal A_j.
\end{equation}
The inclusion holds because every letter of an original upper block
belongs to the original lower dictionary, and shortening only deletes
letters. Thus each column shortening involves at most $g$ blocks (tracks).
At level zero no column shortening is needed, so its support need not have
size at most $g$.

Compression proceeds top-down and renames letters as it goes. Shortening the
level-$(j+1)$ blocks changes which level-$j$ letters occur; those letters are
shortened in turn, so the level-$j$ marker and every lower derived monitor change
as well. We therefore never compare the compressed word with the original
hierarchy directly. Instead we keep, at each level $j$, an injective renaming
$\rho_j$ of the letter support $\mathcal E_j$ under which the new derived monitor
\emph{agrees} with the original one: it has the same transitions, targets, and both
admission relations on the renamed letters, at every state, including the symbolic
states and the start label \eqref{eq:monitor-agreement}. Agreement up to an injective
renaming suffices because a monitor sees letters only through these values and
through equality of letters, and an injective renaming preserves both equality and
inequality. Hence agreement preserves comparisons, runs, and holes. Since the support
contains the marker, injectivity also preserves marker avoidance.
\Cref{lem:simultaneous-compression} composes the agreement given by the column
summaries with agreement at the level below. Each level multiplies expansion lengths
by at most $K$, so a renamed level-$j$ letter expands to at most $K^j$ letters. The
original hierarchy and dictionaries stay unchanged throughout this induction;
intermediate compressed words need not be dictionary-normal.

\paragraph{The Bound}
For a lettered monitor with $N$ states and a tuple of $h$ blocks, put
\begin{equation}\label{eq:col}
  \mathit{col}(N,h) = N^{(N+1)h}\cdot 2^{N(N+1)h}\cdot 2^{h^2}\cdot 2^{(N+1)^2h^2}.
\end{equation}
This product counts all assignments to the summary fields and therefore
bounds the number of reachable summary states.
\begin{center}
\begin{tabular}{@{}ll@{}}
\toprule
Factor & Summary field\\
\midrule
$N^{(N+1)h}$ & Exit states: $h$ blocks, $N+1$ entry states\\
$2^{N(N+1)h}$ & Terminal-comparison bits: also $N$ final-marker states\\
$2^{h^2}$ & Equality bits for ordered pairs of blocks\\
$2^{(N+1)^2h^2}$ & Derived-admission bits for block and entry-state pairs\\
\bottomrule
\end{tabular}
\end{center}
Each shortened block has fewer than $\mathit{col}(N+j,g)\le K$
letters, so with its closing marker it has at most $K$ letters. Expanding down
one level therefore multiplies lengths by at most $K$. The unary top has at most
$N+\Delta$ letters, and there are $\Delta\le d$ expansions. Define the single length bound
\begin{equation}\label{eq:holebound}
 \holeBound(N,d)=(N+d+1)\,
       \mathit{col}\bigl(N+d+1,\DictBound(N,d)\bigr)^d.
\end{equation}
The resulting nonempty hole has depth at most $d$ and length at most
$(N+\Delta)K^\Delta\le\holeBound(N,d)$, proving \cref{thm:search}.
The alphabet size does not enter the bound. \Cref{sec:complexity}
gives a closed-form estimate and distinguishes hole length from
exhaustive enumeration cost.

\FloatBarrier
